\section{The empirical comparison}
\label{sec:comparison}

Two comparisons are available. The securities portfolios can be checked against fair
values the banks disclosed a year after the model was calibrated, which is a test of
the asset-side machinery. And SVB's published economic-value profile can be checked
at the shock points the model was not fitted to, which is a test of the whole
construction including the deposit valuation. Neither comparison was available to
the model: in both cases the figures being compared against were published after the
date the model knows about, and no parameter was adjusted once they were seen.

\subsection{Pricing the portfolios a year later}

The exercise is mechanical. Each position is aged by one year, the curve at
31~December 2022 replaces the curve at 31~December 2021, and the option-adjusted
spreads solved at the earlier date are left where they are.
Table~\ref{tab:oos} reports the result against what the banks disclosed.

\begin{table}[htbp]
\centering
\caption{Out-of-sample securities valuation at 31~December 2022, in price per \$100
of amortised cost net of credit allowance where one is recorded. Spreads were
calibrated at 31~December 2021 and then frozen.
Schwab's held-to-maturity portfolio is omitted because the transfers described in
Section~\ref{sec:objects} make its disclosed cost basis non-comparable; the corrected
comparison appears below.}
\label{tab:oos}
\small
\begin{tabular}{llrrr}
\toprule
Bank & Classification & Predicted & Disclosed & Error (\%) \\
\midrule
SVB    & Held to maturity   & 80.36 & 83.41 & $-3.66$ \\
SVB    & Available for sale & 89.28 & 91.14 & $-2.05$ \\
Schwab & Available for sale & 85.72 & 92.33 & $-7.16$ \\
\bottomrule
\end{tabular}
\end{table}

SVB's held-to-maturity book is the headline case, because it is the portfolio whose
unrecognised loss mattered in March 2023. The model puts it at 80.36 against the
83.41 implied by SVB's disclosure, which is a fair value of \$76,169m against a
carrying amount of \$91,321m -- \$91,327m of amortised cost less a \$6m credit
allowance. The error is $-3.66\%$, and it is in the conservative direction: the
reconstruction understates the value of the book, which is to say it overstates the
loss.

Schwab's figure needs the correction set out in Section~\ref{sec:objects} before it
can be read. Measured against the cost basis Schwab reports, the model is wrong by
$-7.16\%$. Measured against what Schwab originally paid for the securities, which is
the basis on which SVB's portfolio is measured, the comparison is 85.72 predicted
against 87.47, an error of $-1.75\%$. On a whole-portfolio basis the corrected error
is $-2.00\%$, against $-6.35\%$ uncorrected. More than two thirds of the apparent
error in Schwab's case is the accounting election rather than the model.

Taken together with SVB's whole-portfolio error of $-3.34\%$, the pattern worth
noting is not the size of any individual error but the sign of all of them. Every
comparison in Table~\ref{tab:oos}, and every variant in the sensitivity analysis of
Section~\ref{sec:sensitivity}, errs in the same direction. A model that was simply
imprecise would scatter. This one is biased, and
Section~\ref{sec:sensitivity} traces where the bias lives rather than leaving it as
a remark.

\subsection{Reproducing SVB's own disclosed profile}

The second comparison is the more demanding one, because an economic-value figure
depends on the deposit valuation as well as the asset valuation, and the deposit
assumption is not disclosed.

One parameter is solved for. Holding the deposit beta at the 0.60 that SVB's own
filing states, the uninsured-deposit half-life is chosen so that the model returns
the $-13.9\%$ that SVB reported at $+100$ basis points. That solve gives 3.59~years
for uninsured deposits and 12.56~years for insured. Everything else in
Table~\ref{tab:profile} follows from it without further adjustment.

\begin{table}[htbp]
\centering
\caption{Modelled against disclosed change in economic value of equity, SVB at
31~December 2021. The $+100$ basis point row is the calibration target: the deposit
half-life was solved to reproduce it, and the agreement there is arithmetic rather
than evidence. The other rows are out of sample.}
\label{tab:profile}
\small
\begin{tabular}{lrrl}
\toprule
Shock & Modelled (\%) & Disclosed (\%) & \\
\midrule
$+200$bp & $-27.61$ & $-27.70$ & out of sample \\
$+100$bp & $-13.90$ & $-13.90$ & \textit{calibration target} \\
$-100$bp & $4.60$   & $6.00$   & out of sample \\
$-200$bp & $2.54$   & $3.10$   & out of sample \\
\midrule
Base EVE (\$m) & 21,271 & 20,672 & not fitted \\
\bottomrule
\end{tabular}
\end{table}

Three things in that table are worth separating.

The $+200$ basis point result is the strongest single number in the paper. The model
returns $-27.61\%$ against the $-27.70\%$ SVB disclosed, having been fitted only at
$+100$ basis points. Nothing in the construction forces the two to agree at twice
the shock size, because the relationship between shock size and economic value is not
linear in a portfolio containing prepayment options and a floored down-shock.

The base level is not fitted at all. Only the sensitivity was targeted, so the
\$21,271m the model returns against SVB's disclosed \$20,672m is an independent
comparison, and it overstates by 2.9\%. That is a real discrepancy and not a
rounding artefact. It is consistent with the asset-side conservatism in
Table~\ref{tab:oos} being more than offset, at the base date, by a deposit franchise
the model values generously.

The downward shocks are where the qualitative test sits. SVB's disclosure is
non-monotonic: a 100 basis point fall adds 6.0\% to economic value and a 200 basis
point fall adds only 3.1\%, so the second hundred basis points of decline is
value-destroying where the first was value-creating. The model reproduces that shape
independently, falling from $4.60\%$ to $2.54\%$. It understates both levels by
roughly a point and a half, but it turns in the same place and for the same reason:
the post-shock floor absorbs most of the second hundred basis points at a date when
the entire front of the curve sat below 2\%, while the deposit franchise continues to
lose value as the spread it earns compresses. A model without the floor would report
a smoothly improving profile and would miss this entirely.

\subsection{What these comparisons do and do not license}

The asset-side comparison is a genuine out-of-sample test and supports a
correspondingly firm claim: a reconstruction built from public filings and public
curve data priced two large securities portfolios a year forward to within two to
four percent, with a known and consistent directional bias.

The profile comparison supports a narrower claim. It establishes that the framework,
once given one parameter from the bank's own disclosure, reproduces the rest of that
bank's published sensitivity profile including a feature of its shape that is easy to
get wrong. It does not establish that the deposit half-life so obtained is the right
one, because the only evidence for that figure is the disclosure it was fitted to.
That is the subject of Section~\ref{sec:sensitivity}, and it is where the paper's
central result lies.
